\begin{helper}[The path-index-two first-vertex replacement]
Let \((\varphi,T)\) be a minimum bad pair with
\(T=(y_0,e_1,y_1,\ldots,e_n,y_n)\), \(n\geq2\), and \(\rho(T)=2\).
Then there is a tree-sequence \(S\) of the same length whose exact edge
in every position is the corresponding edge of \(T\), whose first two
vertices are \(y_1,y_0\) in that order, and whose later vertices are the
corresponding vertices of \(T\), such that \((\varphi,S)\) is a bad pair
and \(S\) has smaller Tashkinov rank than \(T\).
\end{helper}
\begin{proof}
Write \(T=(y_0,e_1,y_1,\ldots,e_n,y_n)\).  The assertion
\(\rho(T)=2\), together with \noderef{TashTreePathIndex} and
\noderef{TashTreeSuffixPath}, says that every \(e_h\) with
\(3\leq h\leq n\) joins \(y_{h-1}\) to \(y_h\).  The attachment field of
\noderef{TreeSequence} says that \(e_2\) joins \(y_2\) to one of the two
earlier vertices \(y_0,y_1\).

It cannot join \(y_2\) to \(y_1\).  If it did, the first edge
\(e_1=e_0\) joins \(y_0\) to \(y_1\), again by the attachment field (the
only earlier index at the first position is zero), \(e_2\) would join
\(y_1\) to \(y_2\), and all later consecutive equations were just
obtained from the suffix condition.  Thus \(1\) would satisfy
\noderef{TashTreeSuffixPath}.  Leastness in
\noderef{TashTreePathIndex}, together with its lower bound \(1\), would
give \(\rho(T)=1\), contrary to \(\rho(T)=2\).  Hence \(e_2\) joins
\(y_0\) to \(y_2\).

Define
\[
 S=(y_1,e_1,y_0,e_2,y_2,e_3,y_3,\ldots,e_n,y_n).
\]
This is a \noderef{TreeSequence}, as follows field by field.  Its length
is \(n\), and its exact edge in every position is the same edge identity
as in \(T\).  Edge injectivity is therefore inherited.  Its vertex list
is obtained by transposing the distinct entries \(y_0,y_1\), so vertex
injectivity is inherited as well.  Its first edge is still the exact
edge \(e_1=e_0\).  The endpoint pair of that edge is unchanged, because
the pair is unordered.  The preceding paragraph shows that \(e_2\)
joins the new vertex in position two, namely \(y_2\), to the earlier
new vertex \(y_0\) in position one.  Every \(e_h\), \(h\geq3\), joins
the unchanged consecutive vertices \(y_{h-1},y_h\).  These are all the
attachment obligations.  They also establish the entrywise equalities
for \(S\) asserted in the statement.

We verify the Tashkinov condition rather than infer it from notation.
Unfolding \noderef{TashTreeMinimumBadPair} and
\noderef{TashTreeBadPair} supplies \noderef{TashkinovTree} for \(T\).
For \(e_2\), its old missing-color witness is one of \(y_0,y_1\); both
vertices still occur before \(e_2\) in \(S\).  For an edge \(e_h\) with
\(h\geq3\), an old witness at \(y_0\) or \(y_1\) remains earlier after
their transposition, while any other old witness retains its position.
The exact edge and its color are unchanged.  Hence every old Tashkinov
witness yields an earlier witness in \(S\), proving
\noderef{TashkinovTree} for \(S\).

The entrywise vertex description makes
\noderef{TreeSequenceVertexSet} for \(S\) equal to that for \(T\).
Therefore the failure of \noderef{ColoredVertexSetElementary} supplied
by badness of \(T\) also holds for \(S\).  The \(k\)-triple is literally
unchanged.  These three facts give \noderef{TashTreeBadPair} for
\((\varphi,S)\).

Finally the attachment checks above show that the entire displayed
sequence \(S\) is a path in order, so \(1\) satisfies
\noderef{TashTreeSuffixPath}; its required lower bound forces
\(\rho(S)=1\).  Since \(\rho(T)=2\), the first coordinate in
\noderef{TashTreeRankLess} strictly decreases.  Thus \(S\) is the
asserted smaller bad replacement.
\end{proof}
